% ============================
%   5. Implementation
% ============================
\section{Implementation details}
\label{sec:implementation}

\subsection{Technology stack}

\begin{table}[h]
\centering
\small
\begin{tabularx}{\textwidth}{l X}
\toprule
\textbf{Layer} & \textbf{Choice} \\
\midrule
Backend framework            & FastAPI~\cite{ramirez2022fastapi} on Uvicorn \\
ORM                          & SQLAlchemy 2 (async) with aiosqlite \\
Authentication               & OAuth2 password flow $\rightarrow$ JWT signed with HS256 \\
Password hashing             & bcrypt~\cite{bcrypt} (direct calls, no passlib) \\
Packet capture               & Scapy 2.5+~\cite{scapy} \\
Machine learning             & scikit-learn~\cite{pedregosa2011scikit} (RandomForest, IsolationForest) \\
Numerical                    & NumPy \\
Validation                   & Pydantic v2 \\
Frontend                     & React 18 + Vite~\cite{vite} \\
Styling                      & TailwindCSS 3 \\
Charts                       & Recharts \\
Icons                        & lucide-react \\
Test \& packaging            & pytest, joblib, virtualenv \\
\bottomrule
\end{tabularx}
\caption{Technology stack as of the current milestone.}
\end{table}

The choice of FastAPI was driven by three properties: native asynchronous
support (necessary because the packet pipeline is event-loop driven), an
automatically generated OpenAPI schema (useful for the defense
demonstration), and very low boilerplate for dependency-injected
authentication. SQLite was retained for the same reason that motivated
the one-click launcher: zero installation footprint.

\subsection{Packet model}

The internal \code{Packet} dataclass is intentionally framework-agnostic
so that the \code{LiveSniffer} and the \code{TrafficSimulator} can feed
it interchangeably:
\begin{lstlisting}[language=Python,caption={Internal packet representation}]
@dataclass
class Packet:
    ts: datetime            # UTC, timezone-aware
    protocol: str           # "TCP" | "UDP" | "ICMP" | "OTHER"
    src_ip: str
    src_port: int
    dst_ip: str
    dst_port: int
    length: int             # full L2 length
    flags: str              # comma-separated TCP flags
    payload: bytes
\end{lstlisting}

\subsection{Detection pipeline core}

The heart of the system is the \code{Dispatcher.handle} coroutine. It
shows three concerns kept deliberately separate: metrics, publication,
and detection.
\begin{lstlisting}[language=Python,caption={Dispatcher entry point (excerpt)}]
async def handle(self, pkt: Packet) -> None:
    self._pkt_count += 1
    self._ring.append(pkt)
    self._tick_pps()

    await self._publish_traffic(pkt)
    await self._maybe_persist_event(pkt)

    alerts = self._evaluate(pkt)
    if alerts:
        await self._persist_and_broadcast(alerts)
\end{lstlisting}

The merge logic that produces the three alert categories
(signature-only, ml-only, hybrid) lives in \code{\_evaluate}; it is the
locus of the design choice highlighted in Section~\ref{sec:background}.

\subsection{Authentication and authorisation}

Login uses the OAuth2 password flow exposed by FastAPI; tokens are
JWTs (HS256) carrying \code{sub} (username) and an optional role claim,
with a default 60-minute expiry. Every protected route declares
\code{Depends(get\_current\_user)}, which decodes the token, loads the
matching \code{User} row and rejects expired or unknown credentials
with HTTP 401. WebSocket endpoints follow the same path but accept the
token as a query-string parameter, since browser WebSocket clients
cannot set custom request headers~\cite{ietf6455}.

\subsection{Forensic file format}

Each per-alert capture file is a JSON Lines document where every line
encodes one packet (timestamp, 5-tuple, flags, length, base64-encoded
payload preview). The format was chosen over PCAP because:
\begin{itemize}
    \item It is line-delimited, so the dashboard can stream a paginated
        preview without parsing the whole file.
    \item It is human-readable, which matters for academic demonstration.
    \item It avoids the dependency on libpcap-format encoding (Scapy's
        \code{wrpcap}) which complicates portability to operating
        systems without a libpcap install.
\end{itemize}

\subsection{Resilience and recent stabilisation work}

During the latest iteration, several rough edges identified in everyday
use were addressed:
\begin{description}
    \item[First-login race.] On a cold boot the very first
        \code{POST /api/auth/login} could race against the lifespan
        seeding phase and return HTTP 500. The endpoint now catches the
        transient database exception, returns a clean 503, and the
        frontend retries up to three times with linear back-off,
        eliminating the ``mash the button until it works'' artefact.
    \item[Idempotent shutdown.] The launcher's shutdown routine is now
        guarded by a flag so that the trio of atexit, signal handler
        and \code{finally} clause does not run shutdown three times. On
        Windows this previously required the user to confirm the
        ``Terminate batch job?'' prompt twice.
    \item[Authenticated download.] The Forensics page now downloads
        files through a \code{fetch}-based helper that attaches the JWT
        in the Authorization header and converts the response to a
        \code{Blob} URL; the previous \code{<a download>} link failed
        because anchor navigation cannot set custom headers, which
        produced an opaque ``File wasn't available'' Chrome error.
    \item[Filter correctness.] The Alerts page filters (severity,
        method, status, free-text) now apply to both the historical
        rows fetched from the API and to the live WebSocket stream;
        previously the live stream bypassed the filter predicate.
    \item[Themed confirmation dialog.] Destructive actions (currently
        \textit{Delete rule}) raise a navy-cyan modal component
        consistent with the rest of the interface instead of the
        browser's native pop-up.
    \item[ML registry collisions.] Re-training with the default model
        name no longer crashes on the \code{UNIQUE} constraint; the
        backend auto-suffixes a free name and the frontend defaults the
        name per algorithm.
    \item[Settings page.] The previously read-only Settings tab now
        offers a change-password form (backed by
        \code{POST /api/auth/change-password}), surfaces the current
        signed-in user and role, and documents the deployment
        configuration knobs.
\end{description}

\subsection{Source-tree organisation}

\begin{lstlisting}[basicstyle=\ttfamily\scriptsize,caption={Repository layout (abridged)}]
Licenta_NIDS/
|-- backend/
|   |-- app/
|   |   |-- api/                  FastAPI routers
|   |   |-- core/                 config, database, security
|   |   |-- engine/               sniffer, dispatcher, signatures, ML detector
|   |   |-- models/               SQLAlchemy ORM
|   |   |-- schemas/              Pydantic
|   |   |-- ws/                   WebSocket endpoints
|   |   `-- main.py               app factory + lifespan
|   |-- data/
|   |   |-- rules/default_rules.yaml
|   |   `-- models/               *.joblib
|   `-- requirements.txt
|-- frontend/
|   |-- src/
|   |   |-- api/                  HTTP + WS client
|   |   |-- components/           shared UI (modals, badges, sidebar)
|   |   |-- pages/                Login, Dashboard, Alerts, AlertDetail,
|   |   |                         Rules, Forensics, MLModels, Settings
|   |   `-- store/auth.jsx        auth context
|   `-- package.json
|-- docs/progress-report/         this document
|-- run.py                        cross-platform launcher
|-- run.bat                       Windows wrapper
|-- run.sh                        Unix wrapper
`-- README.md
\end{lstlisting}
